\section{단순 대수적 확장}
\label{sec:simple-algebraic-extensions}

\begin{learninggoal}
$K[\alpha]$와 $K(\alpha)$를 구별하고, $\alpha$가 대수적일 때 두 대상이 같음을 증명한다. 최소다항식으로 단순확장의 기저와 연산을 계산한다.
\end{learninggoal}

\subsection{환으로 생성하기와 체로 생성하기}

체확장 $L/K$와 $\alpha\in L$에 대해 다음 두 대상을 구별한다.

\[
\begin{aligned}
  K[\alpha]&=\{f(\alpha):f\in K[X]\},\\
  K(\alpha)&=\text{$K$와 $\alpha$를 포함하는 $L$의 가장 작은 부분체}.
\end{aligned}
\]

첫 번째는 $K$와 $\alpha$가 생성하는 부분환이고, 두 번째는 생성하는 부분체이다. 항상
\[
  K[\alpha]\subseteq K(\alpha)
\]
이지만 일반적으로 등호가 성립하지는 않는다. 예를 들어 $t$가 $K$ 위 초월적이면
\[
  K[t]\subsetneq K(t)
\]
이다. 오른쪽에는 $1/t$가 있지만 왼쪽에는 없다.

\begin{theorem}[대수적 단순확장의 구조]
$\alpha\in L$가 $K$ 위에서 대수적이고 최소다항식을
\[
  m=m_{\alpha,K}
\]
라 하자. 그러면
\[
  K[\alpha]=K(\alpha)
  \cong K[X]/(m).
\]
또한 $n=\deg m$이면
\[
  1,\alpha,\dots,\alpha^{n-1}
\]
가 $K(\alpha)$의 $K$-기저이고
\[
  [K(\alpha):K]=n.
\]
\end{theorem}

\begin{proof}
평가 준동형의 핵은 $(m)$이므로 제1동형정리에 의해
\[
  K[X]/(m)\cong K[\alpha].
\]
$m$은 기약이므로 $(m)$은 극대아이디얼이고 몫환은 체이다. 따라서 $K[\alpha]$ 자체가 체이며, $K$와 $\alpha$를 포함하는 가장 작은 체 $K(\alpha)$와 같아야 한다.

나눗셈 알고리즘에 의해 모든 $f\in K[X]$는
\[
  f=qm+r,
  \qquad
  \deg r<n
\]
으로 유일하게 표현된다. 따라서 모든 원소는
\[
  a_0+a_1\alpha+\cdots+a_{n-1}\alpha^{n-1}
\]
의 꼴로 표현된다. 이 표현이 $0$이면 차수가 $n$보다 작은 다항식 $r$가 $\alpha$를 근으로 갖는다. 최소다항식의 최소성에서 $r=0$이므로 기저가 된다.
\end{proof}

\begin{definition}[단순확장]
어떤 $\alpha\in L$에 대해 $L=K(\alpha)$이면 $L/K$를 \emph{단순확장}(simple extension)이라 한다. 이때 $\alpha$를 원시원 또는 생성원이라 한다.
\end{definition}

\subsection{표준형으로 계산하기}

최소다항식이
\[
  m(X)=X^n+c_{n-1}X^{n-1}+\cdots+c_0
\]
이면 $K(\alpha)$ 안에서
\[
  \alpha^n=-c_{n-1}\alpha^{n-1}-\cdots-c_0
\]
이다. 따라서 높은 거듭제곱을 반복해서 낮추면 모든 계산을 기저
\[
  1,\alpha,\dots,\alpha^{n-1}
\]
의 선형결합으로 정리할 수 있다.

\begin{example}
$\alpha=\sqrt[3]{2}$라 하자. $X^3-2$는 $\Q[X]$에서 기약이므로
\[
  [\Q(\alpha):\Q]=3
\]
이고 모든 원소는
\[
  a+b\alpha+c\alpha^2
\]
로 유일하게 표현된다. 또한 $\alpha^3=2$이므로
\[
  \alpha^5=2\alpha^2.
\]
\end{example}

\begin{example}[역원 계산]
위와 같은 $\alpha=\sqrt[3]{2}$에 대해
\[
  (1+\alpha)(1-\alpha+\alpha^2)=1+\alpha^3=3.
\]
따라서
\[
  (1+\alpha)^{-1}=\frac{1-\alpha+\alpha^2}{3}.
\]
\end{example}

일반적으로 영이 아닌 $f(\alpha)$의 역원은 확장 유클리드 알고리즘으로 구할 수 있다. $m$이 기약이고 $m\nmid f$이면
\[
  \gcd(f,m)=1
\]
이므로 어떤 $a,b\in K[X]$가 존재하여
\[
  af+bm=1.
\]
$X=\alpha$를 대입하면
\[
  a(\alpha)f(\alpha)=1
\]
이다.

\subsection{유한체의 첫 구성}

\begin{example}
$X^2+X+1$은 $\F_2[X]$에서 기약이다. $\alpha$를 그 한 근이라 하면
\[
  \F_2(\alpha)\cong\F_2[X]/(X^2+X+1)
\]
이고
\[
  \alpha^2=\alpha+1.
\]
기저는 $1,\alpha$이며 네 원소
\[
  0,1,\alpha,\alpha+1
\]
를 얻는다. 영이 아닌 원소들은
\[
  \alpha^3=1
\]
을 만족한다.
\end{example}

\begin{keyidea}
최소다항식은 단순확장의 완전한 계산 규칙이다. 그 차수는 확장의 차수이고, 그 방정식은 높은 거듭제곱을 기저 안으로 되돌리는 축약 규칙이다.
\end{keyidea}

\begin{checkpoint}
\begin{enumerate}[label=(\alph*)]
  \item $\Q(i)$의 표준형과 $\Q$-기저를 적어라.
  \item $\F_2[X]/(X^3+X+1)$의 원소 수와 표준형을 구하라.
  \item $\Q(\sqrt2)$에서 $(3+2\sqrt2)^{-1}$을 구하라.
\end{enumerate}
\end{checkpoint}
